\documentclass[answers,12pt,addpoints]{exam} 
\usepackage{import}

\import{C:/Users/prana/OneDrive/Desktop/MathNotes}{style.tex}

% Header
\newcommand{\name}{Pranav Tikkawar}
\newcommand{\course}{01:XXX:XXX}
\newcommand{\assignment}{Homework n}
\author{\name}
\title{\course \ - \assignment}

\begin{document}
\maketitle

\newpage
\section{A General Periodic Fourier--State-Space Spatio--Temporal Model}

\subsection{Model setup}

We consider a real-valued Gaussian spatio--temporal process
\[
  \{X_t(x) : x \in [0,2\pi],\; t \in \mathbb{Z}\},
\]
defined on a one-dimensional spatial domain with periodic boundary conditions (a circle) and discrete time. For each time $t$, the spatial field $X_t(\cdot)$ is represented by a truncated Fourier series
\begin{equation}
  X_t(x)
  \;=\;
  a_0(t)
  \;+\;
  \sum_{j=1}^{J}
    \Bigl[
      a_j(t)\cos(jx) + b_j(t)\sin(jx)
    \Bigr],
  \qquad x \in [0,2\pi],
  \label{eq:Xt-fourier}
\end{equation}
where $J$ is a truncation level and
\[
  a_j(t),\, b_j(t), \quad j=0,\dots,J,\; t\in\mathbb{Z},
\]
are jointly Gaussian random coefficients. The representation \eqref{eq:Xt-fourier} is exact in the limit $J\to\infty$ and is a finite-dimensional approximation in practice.

We collect the coefficients into a state vector
\begin{equation}
  \mathbf{Z}_t
  \;=\;
  \bigl(a_0(t), a_1(t),\dots,a_J(t), b_1(t),\dots,b_J(t)\bigr)^\top
  \;\in\; \mathbb{R}^{2J+1}.
  \label{eq:state-vector}
\end{equation}

\subsection{General linear Gaussian dynamics in Fourier space}

The temporal evolution of the Fourier coefficients is specified by a linear Gaussian state-space model
\begin{equation}
  \mathbf{Z}_{t+1}
  =
  \boldsymbol{\mu}
  +
  M\,(\mathbf{Z}_t - \boldsymbol{\mu})
  +
  \boldsymbol{\varepsilon}_t,
  \qquad
  \boldsymbol{\varepsilon}_t \sim \mathcal{N}(\mathbf{0}, \Sigma_\varepsilon),
  \label{eq:Zt-state-space}
\end{equation}
where
\begin{itemize}
  \item $\boldsymbol{\mu} \in \mathbb{R}^{2J+1}$ is the (possibly nonzero) mean state;
  \item $M \in \mathbb{R}^{(2J+1)\times(2J+1)}$ is the transition matrix;
  \item $\Sigma_\varepsilon \in \mathbb{R}^{(2J+1)\times(2J+1)}$ is the innovation covariance matrix;
  \item $\{\boldsymbol{\varepsilon}_t\}$ are i.i.d.\ Gaussian, independent of $\mathbf{Z}_0$.
\end{itemize}

Equation \eqref{eq:Zt-state-space} implies the process $\mathbf{Z}_t$ is first-order Markov in time. The conditions for stationarity of $\{\mathbf{Z}_t\}$ are standard: if the spectral radius of $M$ is strictly less than $1$, i.e.
\[
  \rho(M) < 1,
\]
then there exists a unique stationary distribution $\mathcal{N}(\boldsymbol{\mu}, \Sigma_0)$ for $\mathbf{Z}_t$, where $\Sigma_0$ solves the discrete-time Lyapunov equation
\begin{equation}
  \Sigma_0
  =
  M\Sigma_0 M^\top + \Sigma_\varepsilon.
  \label{eq:lyapunov}
\end{equation}
Given $\boldsymbol{\mu}$, $M$, and $\Sigma_\varepsilon$, the stationary covariance $\Sigma_0$ is determined by \eqref{eq:lyapunov}. Conversely, one can specify a desired stationary covariance $\Sigma_0$ and transition $M$, and then define
\begin{equation}
  \Sigma_\varepsilon
  =
  \Sigma_0 - M\Sigma_0 M^\top,
  \label{eq:Sigma-eps-from-Sigma0}
\end{equation}
provided the right-hand side is positive semidefinite.

\subsection{Matérn-like spatial margin on the circle}

We seek to construct $\Sigma_0$ such that the induced spatial margin of $X_t(\cdot)$ at each time $t$ is approximately that of a Matérn field on the circle.

In one spatial dimension on $\mathbb{R}$, the Matérn covariance with variance $\sigma^2$, smoothness $\nu > 0$, and range parameter $\ell>0$ can be written as
\[
  C_\nu(h)
  =
  \sigma^2 \frac{2^{1-\nu}}{\Gamma(\nu)}
  \left(
    \sqrt{2\nu}\,\frac{|h|}{\ell}
  \right)^\nu
  K_\nu\!\left(
    \sqrt{2\nu}\,\frac{|h|}{\ell}
  \right),
\]
whose spectral density is proportional to
\[
  S(\omega)
  \;\propto\;
  \bigl(\kappa^2 + \omega^2\bigr)^{-(\nu + 1/2)},
  \qquad
  \kappa = \frac{\sqrt{2\nu}}{\ell}.
\]

On the periodic domain, we approximate this by imposing that the Fourier coefficients at a fixed time have variances
\begin{equation}
  \sigma_j^2
  \;\propto\;
  \bigl(\kappa^2 + j^2\bigr)^{-(\nu + 1/2)},
  \qquad j = 0,\dots,J,
  \label{eq:matern-spectrum-general}
\end{equation}
with $\kappa = \sqrt{2\nu}/\ell$. For simulation, it is convenient to normalize these so that
\begin{equation}
  \sum_{j=0}^{J} \sigma_j^2 = \sigma^2_{\text{tot}},
  \label{eq:sigma-normalization}
\end{equation}
for a chosen total marginal variance $\sigma^2_{\text{tot}}$. The choice of $\nu$ and $\ell$ controls, respectively, the spatial smoothness and range of $X_t(\cdot)$.

A minimal assumption consistent with \eqref{eq:matern-spectrum-general} is that, at stationarity,
\begin{equation}
  \operatorname{Var}(a_j(t)) = \operatorname{Var}(b_j(t)) = \sigma_j^2,
  \qquad j=1,\dots,J,
  \label{eq:marginal-modes}
\end{equation}
and $\operatorname{Var}(a_0(t)) = \sigma_0^2$, with all these variances appearing along the diagonal of $\Sigma_0$ in \eqref{eq:lyapunov}. Cross-frequency and cosine/sine correlations can be incorporated in $\Sigma_0$ without changing the overall shape of the marginal spectral envelope given by \eqref{eq:matern-spectrum-general}.

\subsection{Mode-specific temporal behaviour}

The most general linear Gaussian model allows $\mathbf{Z}_t$ to evolve with a fully populated transition matrix $M$ and arbitrary innovation covariance $\Sigma_\varepsilon$. For purposes of interpretation, it is useful to describe a structured, but still flexible, parameterization.

\paragraph{Frequency-dependent AR(1) persistence.}  
For each frequency $j$, the coefficients may be assigned an effective persistence parameter $\rho_j \in (-1,1)$ that governs their typical temporal autocorrelation length. In the simplest block-diagonal model (no cross-frequency coupling in $M$), one can set
\begin{equation}
  a_j(t+1) = \rho_j a_j(t) + \eta_{a,j,t}, \quad
  b_j(t+1) = \rho_j b_j(t) + \eta_{b,j,t},
  \qquad j=1,\dots,J,
  \label{eq:simple-ar1-modes}
\end{equation}
and
\begin{equation}
  a_0(t+1) = \rho_0 a_0(t) + \eta_{0,t},
  \label{eq:simple-ar1-mean}
\end{equation}
with Gaussian innovations
\[
  \eta_{0,t} \sim \mathcal{N}\bigl(0, \sigma_0^2 (1 - \rho_0^2)\bigr),
  \quad
  \eta_{a,j,t},\eta_{b,j,t} \sim \mathcal{N}\bigl(0, \sigma_j^2 (1 - \rho_j^2)\bigr),
\]
independent across $j$, between $a$ and $b$, and across $t$. Then $\{a_j(t)\}$ and $\{b_j(t)\}$ are stationary AR(1) processes with variance $\sigma_j^2$ and autocorrelation $\rho_j^{|k|}$.

A key nonseparable feature arises from allowing the persistence to depend on frequency, e.g.
\begin{equation}
  \rho_j
  =
  \exp\bigl(-\lambda j^\alpha\bigr),
  \qquad j=1,\dots,J,
  \label{eq:rho-freq-dependence}
\end{equation}
with $\rho_0$ set separately in $(-1,1)$, and parameters $\lambda > 0$ and $\alpha > 0$. In this case, low-frequency modes (small $j$) decorrelate slowly in time, while high-frequency modes decorrelate rapidly, so the effective spatial correlation structure changes with the time lag. This yields a space--time covariance that is stationary but generally nonseparable in space and time.

\paragraph{Cross-frequency and cosine/sine coupling.}  
Beyond \eqref{eq:simple-ar1-modes}--\eqref{eq:simple-ar1-mean}, one can introduce rich dependency structures in either
\begin{itemize}
  \item the transition matrix $M$, allowing the value of mode $j$ at time $t+1$ to depend on multiple modes at time $t$ (e.g.\ banded or low-rank modifications of a block-diagonal structure), or
  \item the innovation covariance $\Sigma_\varepsilon$, allowing simultaneous shocks to be correlated across frequencies and between cosine/sine components.
\end{itemize}
For example, a block-structured $M$ might have $2\times 2$ blocks $M_j$ acting on $(a_j,b_j)$ for each $j$, with
\[
  M_j
  =
  \rho_j
  \begin{pmatrix}
    \cos\theta_j & -\sin\theta_j \\
    \sin\theta_j & \cos\theta_j
  \end{pmatrix},
\]
which induces a rotation (phase drift) and persistence at each frequency. Similarly, $\Sigma_\varepsilon$ could be parameterized as
\[
  \Sigma_\varepsilon = \Lambda \Lambda^\top + D,
\]
with a low-rank factor $\Lambda$ representing a few latent shocks that affect many modes simultaneously, and $D$ diagonal.

\subsection{Space--time covariance}

Under the simple independent-mode AR(1) structure \eqref{eq:simple-ar1-modes}--\eqref{eq:simple-ar1-mean} with zero mean, the space--time covariance
\[
  C(h,k) = \mathbb{E}\bigl[ X_t(x) X_{t+k}(x+h) \bigr]
\]
admits a closed form:
\begin{equation}
  C(h,k)
  =
  \sum_{j=0}^{J}
  \sigma_j^2 \rho_j^{|k|} \cos(jh),
  \label{eq:cov-fourier-ar1-general}
\end{equation}
where $\cos(0\cdot h)=1$ for the $j=0$ term. This follows by substituting \eqref{eq:Xt-fourier}, using the uncorrelatedness of different modes and cosine/sine components, and applying the identity
\[
  \cos A\cos B + \sin A\sin B = \cos(A-B).
\]
Formula \eqref{eq:cov-fourier-ar1-general} makes explicit that the process is stationary in both space and time and that, unless all $\rho_j$ are equal, $C(h,k)$ cannot be written as a separable product $C_{\text{space}}(h) C_{\text{time}}(k)$.

When $M$ and $\Sigma_\varepsilon$ are fully general, the covariance structure can still be characterized via the stationary covariance $\Sigma_0$ from \eqref{eq:lyapunov}, but the expression for $C(h,k)$ in physical space involves both frequency--lag-dependent auto- and cross-covariances among the Fourier coefficients and generally cannot be reduced to \eqref{eq:cov-fourier-ar1-general}.

\subsection{Parameter overview}

The full model includes the following parameters, with varying levels of restriction depending on the chosen structure:

\begin{itemize}
  \item \textbf{Fourier truncation:} $J$ (integer, $J\ge 0$). Controls the finest spatial scale represented.

  \item \textbf{Matérn spatial parameters:}
    \begin{itemize}
      \item $\nu > 0$ (smoothness): larger $\nu$ gives smoother spatial fields.
      \item $\ell > 0$ (range): larger $\ell$ increases the spatial correlation length.
      \item $\sigma^2_{\text{tot}} > 0$ (total marginal variance): overall scale.
      \item $\kappa = \sqrt{2\nu}/\ell$ (inverse range): internal scale parameter used in the spectral envelope \eqref{eq:matern-spectrum-general}.
    \end{itemize}

  \item \textbf{Temporal persistence parameters:}
    \begin{itemize}
      \item $\rho_0 \in (-1,1)$: AR(1) persistence for the mean mode $a_0(t)$.
      \item $\rho_j \in (-1,1)$, $j=1,\dots,J$, e.g.\ via \eqref{eq:rho-freq-dependence} with
        \begin{itemize}
          \item $\lambda > 0$: controls overall decay speed across frequencies,
          \item $\alpha > 0$: controls how strongly persistence changes with $j$.
        \end{itemize}
    \end{itemize}

  \item \textbf{State-space parameters:}
    \begin{itemize}
      \item $\boldsymbol{\mu} \in \mathbb{R}^{2J+1}$: mean state (for zero-mean processes, $\boldsymbol{\mu} = 0$).
      \item $M \in \mathbb{R}^{(2J+1)\times(2J+1)}$: transition matrix. Possible structures:
        \begin{itemize}
          \item block-diagonal in frequencies (independent modes),
          \item block-diagonal with $2\times 2$ rotations per mode,
          \item banded, low-rank, or fully general.
        \end{itemize}
      \item $\Sigma_\varepsilon \in \mathbb{R}^{(2J+1)\times(2J+1)}$: innovation covariance. Possible structures:
        \begin{itemize}
          \item diagonal (independent innovations),
          \item diagonal plus low-rank (factor structure),
          \item banded or fully general symmetric positive semidefinite.
        \end{itemize}
    \end{itemize}

  \item \textbf{Time horizon:} $T$ (integer, $T\ge 1$). Number of time points simulated or observed.

\end{itemize}

By choosing simple structures (e.g.\ block-diagonal $M$ with $\rho_j$ given by \eqref{eq:rho-freq-dependence}, diagonal $\Sigma_\varepsilon$ determined by \eqref{eq:Sigma-eps-from-Sigma0}, and $\boldsymbol{\mu}=0$) one obtains a tractable, interpretable nonseparable model. More complex structures in $M$ and $\Sigma_\varepsilon$ can then be introduced gradually to capture richer cross-frequency interactions while preserving the Matérn-like spatial margin and Markovian temporal dynamics.

\end{document}